\section{Discussion}\label{sec:discussion}

\subsection{The norm, not the labels}
Both comparisons changed the norm of the supervised loss and nothing else. Both lowered the
stress error, on the validation and on the held-out instances: in three dimensions the
median von Mises error by a factor of \numNlThreeNormVm{}, in two the mean by more than half
(\cref{sec:norm}). In two dimensions, measured after the run, the errors of the two
transformers were of about the same size and lay at different places in the spectrum.
Removing the solutions with the energy norm kept raised the median energy-norm and von Mises
errors by factors of at most \numNlThreeLabelWithin{} (\cref{sec:labelfree}), in a comparison
that also changes the weighting of the load cases and the learning rate. The label-free
transformer's advantage over displacement supervision thus lay in the norm of its loss rather
than in the absence of labels. By \cref{cor:training,prop:stress} the energy objective is, up
to a constant, the compliance-weighted mean-square error of the stress field, which is what
the displacement error under-weights. With labels at hand, training in the energy norm with
them was the more accurate in three dimensions and showed no difference beyond the guard in two. The label-free objective is thus not necessarily
the most accurate way to train in the
energy norm, but it is the way to train in it without solutions. The energy term of the
displacement-supervised transformer, which adds the energy to the displacement loss, is a step towards the energy norm. It removed most of the large failures, and the label-free
transformer, which has only that term, was more accurate still (\cref{tab:3d}).

\subsection{Engineering use}\label{sec:discussion:engineering}
For surrogates used in stress analysis, such as those of the families of
\cref{sec:methods:data}, the results suggest three rules. First, train in the energy norm,
which by \cref{cor:training,prop:stress} is, up to a constant, the compliance-weighted error
of the stresses, so a surrogate trained in it is in effect supervised on its stresses.
Given the assembled system, the energy norm of an error costs one sparse matrix-vector
product per load case. A supervised loss can therefore be measured in it with the labels
already computed, on which it depends only through the product $F^{\top}\Ustar$ of each load case (\cref{sec:methods:training}). Where no solutions have been computed, minimising the
energy trains in the same norm without them. Second, compare surrogates by their energy-norm and stress errors, not by their displacement
errors alone. Surrogates compared by displacement error alone can be ranked opposite to their
stresses (\cref{sec:checks:ranking}), and an evaluation can report the energy norm beside the
displacement error. Third, screen
predictions with the energy before they are used. Its sign flags gross failures, those worse
than the zero field, without a solution (\cref{sec:checks:zero}). Each of the \numHoldTwoWorseInstKnorm{}
predictions of the energy-norm-supervised transformer with a mean
relative energy gap above one on the two-dimensional held-out instances (\cref{sec:holdout})
had a load case that the screen flags. Differences of energy rank the predictions of a load case
exactly (\cref{cor:rank}), so, of several predictions from different networks or seeds
or from a network and a Ritz projection, the one of lowest energy is the most accurate in the
energy norm. A prediction that fails should be sent back to the solver. Rescaling it by
$\cstar$ (\cref{prop:amp}) makes it no worse than the zero field without a solve, but not
necessarily accurate.

The gain is in the accuracy of the stresses, not in speed. Training in the energy norm leaves
the cost of a prediction unchanged. Conjugate gradients given the stiffness matrix, stopped by
a test that uses the reference solution, reached the label-free transformer's accuracy within a few
times its inference time (\cref{sec:cost:inference}). Training without labels saves only their
solves, which matters where those are expensive (\cref{sec:cost:labels}). Surrogates, and with them these
rules, are worth their training only where one family of components is analysed many times.
In three dimensions the training of one network was recovered against the direct solve after
about \numCostBreakEvenDirect{} predicted instances (\cref{sec:cost:labels}), with the
caveats of \cref{sec:cost:inference}, under which a faster solver would raise that number. In two dimensions the direct solve took \numCostTwoDirectRatio{} times as long as the inference
and is exact, so a surrogate of plates of this size saves little or no time. The screen needs the stiffness matrix and the load vectors of the instance. Their assembly
took a median of \numCostAsmVal{} against \numCostSurrVal{} for the inference in three
dimensions (\cref{tab:cost}). It is also most of the cost of a Ritz projection, which took a further
\numCostRitzValMin{} to \numCostRitzValMax{} and, at the top of its ladder, was more accurate
than the label-free transformer, in the seed means, on every metric but peak stress
(\cref{sec:cost:ritz}). Passing the screen does not make a prediction accurate
(\cref{sec:checks:zero}), and the accuracy of a surrogate on a family has to be established on
solved instances of that family, as the validation instances establish it here.

In the assessment of parts whose internal voids differ from one part to the next, each part is a new instance of a family. The solids here, with at most three spherical cavities under
four load cases, are a stylised version of that setting. Real voids are more numerous, smaller and irregular. Resolving them needs meshes finer than any on which these networks
were accurate (\cref{sec:discussion:transfer}), and other geometries, numbers of voids and
loads were not tested. Label-free training would be most useful where each solve is expensive, as
in the nonlinear and time-dependent problems of \cref{sec:discussion:outlook}.

As in an energy-based model
\cite{LeCun2007EBM}, which uses one scalar energy to find the most compatible output, to rank
candidates and to detect incompatible ones, the potential energy here is used to train, to rank and to check. It is assembled rather than learned, and exact. The identity does not
depend on the approximation: the energy evaluates a Ritz solution as it does a network's prediction. The rescaling by $\cstar$, the projection onto the span of a prediction
(\cref{prop:amp}), is the smallest enrichment of a finite-element space by it; richer
enrichments have been analysed \cite{Barucq2026Enrich}. Where the stiffness matrix is not at
hand, as for solutions exported from another code, a loss on the strains or the stresses is
the natural substitute; we did not test one.

\subsection{Earlier comparisons}\label{sec:discussion:earlier}
The account is consistent with part of the earlier reports. In soft-tissue beams, adding the
deformation gradient to the supervised loss, a step towards a norm of the strains, narrowed
the gap between supervised and label-free training \cite{Dalton2023}. For a scalar multiscale
problem an energy-norm loss gave lower errors of the solution than $L^2$,
relative $L^2$ or $L^1$ losses on the same labels \cite{Haltmayer2026LOD}. One comparison,
reported in a preprint, points the other way. In the pressure vessels of Sun
et~al.~\cite{Sun2026FEVessel}, the same point-cloud transformer trained on the relative
error of the displacement labels instead of the energy, with every other setting unchanged,
was more accurate in displacement and in the $H^1$ seminorm (their Table 2), against the
three-dimensional result of \cref{sec:labelfree:three}. The two studies differ in ways that
may matter. That comparison is a single run, and its reported errors are the lowest test
errors over the last third of training. Its reference solutions use reduced-integration
hexahedra with hourglass control, and its energy loss takes derivatives from the shape
functions of those elements, with a quadrature that is not stated. \Cref{lem:exact} makes
energy training exact regression on the reference solution only when the energy minimised is
the assembled energy of the discretisation that defines the reference. If the two differ,
the minimiser of the energy loss is not the reference solution, and energy training
regresses towards a different target. Whether this explains the reversal is not known. The
comparison shows that the advantage of label-free training over displacement supervision
reported here is not universal.

\subsection{Transfer}\label{sec:discussion:transfer}
No network transferred to the finer mesh. The readings of supplementary
Section~\SuppTransfer{} point to the output scaling. The output is multiplied by the largest nodal load, which depends
on the mesh while the solution does not. The networks of the July run, whose output was not scaled in this way, lost little accuracy on meshes
\numCoarsenLow{} and \numCoarsenHigh{} times finer than those they were trained on, though
that evidence is weak. Until an output scale independent of the mesh is shown to remove the
failure, the label-free transformer of this paper should be regarded as valid within the range of mesh sizes it was trained on. The discretisation invariance reported for
point-cloud operators trained on the energy \cite{Wang2026PretrainFEM,Sun2026FEVessel} was
tested by varying the density of the input point cloud rather than by refining the mesh.
Label-free fine-tuning on the finer mesh, which needs no solves, was not tested here.

\subsection{Limitations}
The study is restricted to linear elastostatics, where \cref{lem:exact} is exact; one family
of geometries and one set of four load cases per dimension; one model size; and three seeds
per network, for which the noise guard is coarse. \Cref{prop:stress} assumes a homogeneous
isotropic material and an exactly integrated stiffness, as with the linear elements used here. The reported von Mises error weights elements equally, so the proposition
explains it rather than bounds it, although the area- and volume-weighted errors, measured
after the runs, differed little. None of the metrics isolates the stresses at the holes or cavities,
and where the peak of the discrete solution lies at the clamped edge, the peak stress
error does not measure the concentration at a hole (\cref{sec:methods:metrics}). Only the transformer was trained in the energy norm. The graph network was trained on the displacement error alone, and whether its rough stress
fields come from the norm or from message passing over the mesh edges was not tested. The
learning rates were fixed before the runs and not tuned, and the energy-norm-supervised
transformer used the displacement-supervised settings unchanged. The departures from the
pre-registrations are stated with their results: the statistic of the three-dimensional
energy-norm run (\cref{sec:norm:three}), the registration of the two-dimensional comparison
(\cref{sec:norm:two}), the amended three-dimensional gate (\cref{sec:labelfree:three}) and the claim that did not replicate (\cref{sec:holdout}). The spectral readings and the Ritz
baseline were added after the runs. A pre-registered comparison kept the joint-embedding term, a self-supervised term in which
the network predicts the latent vectors of masked regions of the mesh. The scope of this paper, the energy
objective without that term, was chosen after that comparison (supplementary
Section~\SuppPrereg{}). The cost of full self-attention grows with the square of
the number of nodes, which limits the size of the meshes that the transformer can take.

\subsection{Outlook}\label{sec:discussion:outlook}
The construction extends wherever each step of a computation is a convex minimisation. In
linear dynamics, every implicit step of a Newmark-type integrator solves a symmetric
positive definite system whose effective stiffness combines the stiffness, mass and damping
matrices, and \cref{lem:exact} holds step by step \cite{CaoSong2026}. Incremental
variational formulations of dissipative materials write each step of a nonlinear history as
the minimisation of an incremental potential \cite{OrtizStainier1999}, and variational
integrators derive each step of a time integration from a discrete variational principle
\cite{MarsdenWest2001}. Where such a step is a minimisation, its potential remains a
label-free training objective, although the energy gap is no longer a squared norm, and the
counterparts of \cref{prop:amp,cor:zero} remain to be worked out. At finite strains the incremental potential is in general not convex \cite{OrtizStainier1999}, so its
minimiser need not be unique. These are the problems in which a solve may be expensive
enough for label-free training to recover its cost, and in which the choice of norm may again
decide what a surrogate learns.
